\chapter{Lista 8}

\section{Exercício 1, 2 e 3}

Vide documento escaneado PDF ao final do arquivo.

\section{Exercício 4}

Todo o código usado para este exercício está no Anexo \ref{cap:ex4-codigo}.

\subsection{Nós equidistantes no triângulo de referência}

Plotando apenas pontos 
de forma equidistante
no triângulo de referência para $N = 3$, $N = 5$ e $N = 8$, obtemos o resultado exibido
na Figura \ref{fig:ex4-pontos-equidist}.

\begin{figure}[htbp]
	\caption{\label{fig:ex4-pontos-equidist}Nós equidistantes no triângulo de referência para diferentes ordens polinomiais.}
	\begin{center}
        \includegraphics[width=\textwidth,trim=1 1 1 1,clip]{ex4-pontos-equidist.png}
	\end{center}
\end{figure}

Vemos que os nós seguem as regras que vimos em sala de aula:

\begin{itemize}
    \item Número de vértices: fixo em 3
    \item Número de nós nas arestas (excluindo os vértices): $3(N-1)$
    \item Número de nós internos : $\frac{(N-1)(N-2)}{2}$
\end{itemize}

Calculando o número de condição da matriz de Vandermonde para os mesmos valores de $N$,
obtendo a Tabela \ref{tab:cond-vandermonde-equidist}

\begin{table}[h!]
	\caption{Número de condição da matriz de Vandermonde para diferentes valores de $N$ equidistantes no triângulo de referência.}
	\vspace{0.5cm}
	\centering
		\begin{tabular}{c|c}

		$N$ & Número de Condição \\ \hline
		$3$    & $312.4$      \\ \hline
		$5$    & $39510$    \\ \hline
		$8$    & $7.06 \cdot 10^7$     \\
		\end{tabular}
	\label{tab:cond-vandermonde-equidist}
\end{table}

\subsection{Nós via Warp and Blend}

Agora plotamos os pontos via Warp and Blend no triângulo de 
referência para $N = 3$, $N = 5$ e $N = 8$, obtendo o resultado exibido
na Figura \ref{fig:ex4-pontos-warp-blend}, e com o número de condição da matriz de Vandermonde para os mesmos valores de $N$,
obtendo a Tabela \ref{tab:cond-vandermonde-warp-blend}.

\begin{figure}[htbp]
	\caption{\label{fig:ex4-pontos-warp-blend}Nós via Warp and Blend no triângulo de referência para diferentes ordens polinomiais.}
	\begin{center}
        \includegraphics[width=\textwidth,trim=1 1 1 1,clip]{ex4-pontos-warp-blend.png}
	\end{center}
\end{figure}

\begin{table}[h!]
	\caption{Número de condição da matriz de Vandermonde para diferentes valores de $N$ via Warp and Blend no triângulo de referência.}
	\vspace{0.5cm}
	\centering
		\begin{tabular}{c|c}

		$N$ & Número de Condição \\ \hline
		$3$    & $8.21$      \\ \hline
		$5$    & $38.75$    \\ \hline
		$8$    & $400$     \\
		\end{tabular}
	\label{tab:cond-vandermonde-warp-blend}
\end{table}

Verificamos que o número de condição da matriz de Vandermonde para os nós via Warp and Blend 
são significativamente menores do que para os nós equidistantes para uma mesma ordem polonimial $N$,
o que torna o método mais estável para o DGTD 2D do que usar os nós equidistantes. 

\section{Exercício 5}

Vide documento escaneado PDF ao final do arquivo.

\section{Exercício 6}

Todo o código usado para este exercício está no Anexo \ref{cap:ex6-codigo}.

\subsection{Solução Analítica}

Primeira coisa é calcular o gradiente de \(A_z\):

\[
A_z(x,y)=x^2+y^2,
\]

\[
\nabla A_z
=
\frac{\partial A_z}{\partial x}\hat{\mathbf x}
+
\frac{\partial A_z}{\partial y}\hat{\mathbf y}
=
2x\,\hat{\mathbf x}
+
2y\,\hat{\mathbf y}.
\]

Produto vetorial com \(\hat{\mathbf z}\):

\[
\begin{aligned}
\mathbf{B}
&=
\left(2x\,\hat{\mathbf x}
+
2y\,\hat{\mathbf y}\right)
\times
\hat{\mathbf z} \\
&=
2x\left(\hat{\mathbf x}\times\hat{\mathbf z}\right)
+
2y\left(\hat{\mathbf y}\times\hat{\mathbf z}\right).
\end{aligned}
\]

Utilizando as identidades

\[
\hat{\mathbf x}\times\hat{\mathbf z}
=
-\hat{\mathbf y},
\qquad
\hat{\mathbf y}\times\hat{\mathbf z}
=
\hat{\mathbf x},
\]

obtém-se

\[
\begin{aligned}
\mathbf{B}
&=
-2x\,\hat{\mathbf y}
+
2y\,\hat{\mathbf x} \\
&=
2y\,\hat{\mathbf x}
-
2x\,\hat{\mathbf y}.
\end{aligned}
\]

Portanto,

\[
\boxed{
\mathbf{B}(x,y)
=
\begin{bmatrix}
2y\\
-2x\\
0
\end{bmatrix}.
}
\]

\subsection{Solução Numérica}

Definimos $N$ como o número de nós da malha. A solução obtida, calculada
com a função \verb|np.gradient| plotada com a função
\verb|quiver| para $N =3$ e $N = 4$ está exibida na Figura \ref{fig:ex6-campoB}.

\begin{figure}[htbp]
	\caption{\label{fig:ex6-campoB}Solução numérica obtida para o campo $B$ para $N = 3$ e $N = 4$.}
	\begin{center}
        \includegraphics[width=\textwidth,trim=1 1 1 1,clip]{ex6-campoB.png}
	\end{center}
\end{figure}
 
\section{Exercício 7, 8 e 9}

\textcolor{red}{Devido à dificuldades com a implementação do DGTD 2D,
não consegui implementar o método para os exercícios 7, 8 e 9 a tempo antes da data-limite de entrega
da lista.}